\chapter*{术语表}
\phantomsection
\addcontentsline{toc}{chapter}{术语表}

\begin{description}
  \item[可执行文件证据] 从二进制文件、工具输出、测试和报告中得到的可复查事实，例如 checksum、machine、entry、section 和 symbol。
  \item[ELF header] ELF 文件开头的固定结构，保存 class、endianness、machine、entry 和 section/program header 位置等信息。
  \item[section header table] section 目录表，记录每个 section 的名字偏移、类型、flags、地址、文件偏移和大小。
  \item[string table] 字符串表。ELF 结构里很多名字通过偏移引用它，而不是直接保存字符串。
  \item[symbol table] 符号表，把函数或对象名字、绑定、类型、section index、地址和大小连接起来。
  \item[diagnostic] 诊断信息。解析器遇到不支持或不安全的输入时，用 diagnostic 说明原因，而不是继续读垃圾。
  \item[benchmark smoke] 小型 benchmark 检查。它证明 benchmark 入口能跑，但不等于完整性能结论。
  \item[交叉验证] 用两个独立来源检查同一事实，例如用本卷解析器和 readelf 同时验证 machine 或 section size。
\end{description}
